\documentclass[11pt, a4paper]{article}

\usepackage[T1]{fontenc}
\usepackage[utf8]{inputenc}
\usepackage{tgpagella}
\usepackage[spanish, es-noshorthands]{babel}
\usepackage{amsmath, amssymb}
\usepackage[margin=2.5cm]{geometry}
\usepackage{fancyhdr}

\pagestyle{fancy}
\fancyhf{}
\setlength{\headheight}{16pt}
\lhead{\textbf{UCB Sede Tarija} | Ing. Mecatrónica}
\rhead{IMT-121 | Notas del Instructor S06-S09}
\renewcommand{\headrulewidth}{0.4pt}

\begin{document}

\section*{Notas de Planificación Semanal: Unidad 2 (S06)}

\subsection*{Objetivo Estratégico de la Sesión S06-1}
Esta sesión marca el inicio formal de la Unidad 2 tras el Hito 1. La prioridad es consolidar la noción de \textbf{Punto de Operación (DC)} antes de intentar explicar modelos dinámicos ($r_\pi$, $g_m$). El estudiante no debe pensar en AC hasta que entienda de dónde salen los valores estáticos.

\subsection*{Por qué se usa el Circuito de Polarización Fija}
El circuito BJT con emisor a tierra y divisor de base fijo se utiliza deliberadamente por su \textbf{transparencia algebraica}.
\begin{itemize}
    \item Hace muy evidente la suposición de región activa y el cálculo de $V_{CE}$.
    \item Hace flagrante el problema de la sensibilidad a $\beta$ ($I_C \propto \beta$).
    \item \textbf{No} se presenta como un buen diseño de amplificador (esto motiva las topologías estabilizadas posteriores).
\end{itemize}

\subsection*{Esquemas Explícitos en Ejercicios}
Como recomendación pedagógica, los estudiantes no deben verse obligados a reconstruir la topología del circuito a partir de texto durante esta primera sesión. Ya están enfrentando la carga cognitiva de (1) la topología del BJT, (2) la suposición de estado, y (3) la ecuación de la recta de carga. Proveer los esquemas completos elimina una barrera innecesaria y asegura que evaluamos razonamiento en polarización DC.

\end{document}
